\section{An explicit finite-volume strong-coupling bound}
\label{sc:section}

This section gives a quantitative version of stability of a gapped product
Hamiltonian under bounded local perturbations. The argument follows the
Kirkwood--Thomas method used by Bravyi, DiVincenzo, and Loss
\cite{BDL}, with Hilbert-space norms on excitation sectors in place of
absolute values of qubit coefficients. The existence of a sufficiently
small perturbative window, including for infinite-dimensional site spaces,
is established in the literature \cite{Yarotsky}. The purpose here is to
track constants for the specified Hamiltonian lattice gauge model.
Neither the method nor the existence of strong-coupling lattice gaps is
claimed to be new. The derivation and associated records originated in
the research archive \cite{Archive}; the earlier spectral-bounds paper
\cite{Spectral} gives the corresponding non-explicit stability application.

\subsection{A finite-system stability criterion}

Let $\mathcal L$ be a nonempty finite set of sites. At each site $u$, let
$\mathcal H_u$ be a separable Hilbert space and let $h_u$ be a nonnegative
self-adjoint operator with compact resolvent. Assume that $h_u$ has a
normalized simple zero eigenvector $\Omega_u$ and that, for a common
$\Delta>0$,
\begin{equation}
 h_u\big|_{\Omega_u^\perp}\ge\Delta.
 \label{sc:onsite-gap}
\end{equation}
The inequality is vacuous if $\mathcal H_u$ is one-dimensional. On
$\mathcal H=\bigotimes_{u\in\mathcal L}\mathcal H_u$, set
\begin{equation}
 H_0=\sum_{u\in\mathcal L}h_u,\qquad
 V=\sum_{p\in\mathcal P}V_p,\qquad
 H(\varepsilon)=H_0+\varepsilon V.
 \label{sc:general-H}
\end{equation}
Here $\mathcal P$ is a finite indexed family of nonempty site sets, each
of cardinality at most an integer $s\ge1$. Each $V_p$ is self-adjoint, supported on
$p$, and has norm at most $J$. Every site belongs to at most $D$
members of $\mathcal P$, counted with their multiplicities. We may take
$D,J>0$; when there is no perturbation the conclusion below is immediate.
The coupling $\varepsilon$ is real.

Define the dimensionless parameter and positive constants
\begin{align}
 \lambda&=\frac{|\varepsilon|DJ}{\Delta},
 \label{sc:lambda}\\
 \gamma_0(s)^2&=\sum_{t=1}^{s}\frac{\binom{s-1}{t-1}}{t^2},
 &\gamma_1(s)^2&=(s+1)^2\sum_{t=0}^{s}
                        \frac{\binom{s}{t}}{(t+1)^2},
 \label{sc:gammas}\\
 \beta_0(s)&=\gamma_0(s),
 &\beta_k(s)&=2^k\bigl(\gamma_0(s)s^k+
                         k\gamma_1(s)s^{k-1}\bigr)
                         \quad(1\le k\le2s),
 \label{sc:betas}\\
 q_s(r)&=\sum_{k=0}^{2s}\frac{\beta_k(s)}{k!}r^k.
 \label{sc:majorant}
\end{align}

\begin{theorem}[Dimension-independent finite-volume stability]
\label{sc:stability}
Suppose there is a number $r>0$ such that
\begin{equation}
 \lambda q_s(r)\le r,
 \qquad 2\lambda q_s'(r)<1.
 \label{sc:criterion}
\end{equation}
Then $H(\varepsilon)$ has a simple ground eigenvalue $E_0$ and
\begin{equation}
 \operatorname{spec}H(\varepsilon)
 \subset\{E_0\}\cup[E_0+\Delta/2,\infty).
 \label{sc:gap}
\end{equation}
The criterion depends on $s,D,J,\Delta$ and $\varepsilon$, but not on
the number of sites, their Hilbert-space dimensions, or the upper
onsite spectrum.
\end{theorem}

We first prove the theorem when all site spaces are finite-dimensional.
The compact-resolvent extension is proved at the end of this subsection.

For $M\subset\mathcal L$, let
$\mathcal K_M=\bigotimes_{u\in M}\Omega_u^\perp$ and
$\Omega_M=\bigotimes_{u\in M}\Omega_u$. The full Hilbert space is the
orthogonal sum of the sectors
$\mathcal K_M\otimes\Omega_{\mathcal L\setminus M}$.
We identify a sector with $\mathcal K_M$, and write $\Pi_M$ for the
corresponding projection followed by this identification. In particular,
\begin{equation}
 \sum_{M\subset\mathcal L}\|\Pi_M\xi\|^2=\|\xi\|^2.
 \label{sc:orthogonal-sectors}
\end{equation}
Let $\Omega=\bigotimes_{u\in\mathcal L}\Omega_u$.
For $c_M\in\mathcal K_M$, $M\ne\varnothing$, define
\begin{equation}
 A[c_M]=|c_M\rangle\langle\Omega_M|\otimes
                  I_{\mathcal L\setminus M},
 \qquad
 C=\sum_{M\ne\varnothing}A[c_M],
 \qquad
 \|C\|_1=\max_{u\in\mathcal L}
                 \sum_{M\ni u}\|c_M\|.
 \label{sc:creation-norm}
\end{equation}
The decomposition of $C$ is unique, as is seen by applying it to $\Omega$.
We call $C$ a creation operator. The coefficients $c_M$ are vectors,
including when the excitations within a sector are entangled.

Two such elementary operators commute. If their supports intersect,
their product in either order is zero: a vacuum bra on an overlapping
site annihilates an excited factor. If their supports are disjoint,
they commute by tensor factorization. Consequently all creation operators
commute, and every product of more than $|\mathcal L|$ elementary creation
operators vanishes. Their exponentials are therefore finite polynomials.
Also, $\|A[c_M]\|=\|c_M\|$.

Put $H_0^M=\sum_{u\in M}h_u$ on $\mathcal K_M$. Then
\begin{equation}
 H_0^M\ge\Delta|M|,
 \qquad [H_0,A[c_M]]=A[H_0^M c_M].
 \label{sc:free-commutator}
\end{equation}
Thus a second commutator of $H_0$ with creation operators is zero.
We use $\operatorname{ad}_C(X)=[C,X]$.

\begin{lemma}[Support and commutator depth]
\label{sc:support}
Fix $p\in\mathcal P$ and elementary creation operators
$A[c_{M_1}],\ldots,A[c_{M_k}]$. Set $N_*=\bigcup_{i=1}^kM_i$ and
\[
 X=\operatorname{ad}_{A[c_{M_1}]}\cdots
                   \operatorname{ad}_{A[c_{M_k}]}(V_p).
\]
If $X\ne0$, each $M_i$ intersects $p$. If $\Pi_MX\Omega\ne0$, then
\begin{equation}
 N_*\setminus p\subset M\subset N_*\cup p,
 \qquad M=(N_*\setminus p)\cup T,
 \quad T\subset p.
 \label{sc:sector-support}
\end{equation}
Moreover, for arbitrary creation operators $C_1,\ldots,C_k$,
\begin{equation}
 \operatorname{ad}_{C_1}\cdots\operatorname{ad}_{C_k}(V_p)=0
 \quad\hbox{if }k>2s.
 \label{sc:depth}
\end{equation}
\end{lemma}
\begin{proof}
The commutator maps of commuting creation operators commute. A factor
whose support misses $p$ can therefore be placed next to $V_p$ and
annihilates it. Expand a surviving nested commutator into products with
some elementary operators to the left of $V_p$ and the rest to its right.
Outside $p$, the operator $V_p$ changes nothing; every site created there
must remain excited, and a site created twice makes that product zero.
This proves \eqref{sc:sector-support}. In a nonzero product the creation
supports on each side of $V_p$ are pairwise disjoint. Since each meets
$p$, at most $|p|\le s$ can occur on either side. More than $2s$
factors therefore give zero.
\end{proof}

\begin{lemma}[Sector norm estimate]
\label{sc:sector-estimate}
For $0\le k\le2s$ and creation operators $C_1,\ldots,C_k$, define
$\mathcal B_k(C_1,\ldots,C_k)$ by its coefficients
\begin{equation}
 b_M=(H_0^M)^{-1}\Pi_M
       \operatorname{ad}_{C_1}\cdots\operatorname{ad}_{C_k}(V)\Omega,
 \qquad M\ne\varnothing.
 \label{sc:Bk}
\end{equation}
For $k=0$ the commutator product is the identity. Then
\begin{equation}
 \|\mathcal B_k(C_1,\ldots,C_k)\|_1
 \le\frac{DJ}{\Delta}\beta_k(s)
                      \prod_{i=1}^k\|C_i\|_1.
 \label{sc:multilinear}
\end{equation}
\end{lemma}
\begin{proof}
Expand each $C_i$ into its elementary terms, and fix a tuple
$(M_1,\ldots,M_k,p)$. Denote its commutator by $X$ as above. Its norm
satisfies
\begin{equation}
 \|X\Omega\|\le2^kJ\prod_{i=1}^k\|c^{(i)}_{M_i}\|.
 \label{sc:tuple-norm}
\end{equation}
Only tuples with $M_i\cap p\ne\varnothing$ need be retained.
Fix a site $u$ when estimating the norm in
\eqref{sc:creation-norm}. We split the sum into tuples with $u\in p$
and those with $u\notin p$. In the latter case, a nonzero output
sector containing $u$ requires $u\in M_j$ for at least one $j$;
summing over all such $j$ is an upper bound.

If $u\in p$, put $n=|N_*\setminus p|$. A contributing output
sector has the form \eqref{sc:sector-support} with $u\in T$ and
$|M|=n+|T|$. Cauchy--Schwarz and
\eqref{sc:orthogonal-sectors} give
\begin{align}
 \sum_{M\ni u}\frac{\|\Pi_MX\Omega\|}{|M|}
 &\le\left(\sum_{\substack{T\subset p\\u\in T}}
                     \frac{1}{(n+|T|)^2}\right)^{1/2}\|X\Omega\|
 \nonumber\\
 &\le\gamma_0(s)\|X\Omega\|.
 \label{sc:type-zero}
\end{align}
For each fixed $p$ and $i$,
\begin{equation}
 \sum_{M_i:M_i\cap p\ne\varnothing}\|c^{(i)}_{M_i}\|
 \le\sum_{v\in p}\sum_{M_i\ni v}\|c^{(i)}_{M_i}\|
 \le s\|C_i\|_1.
 \label{sc:incidence-sum}
\end{equation}
There are at most $D$ choices of $p$ containing $u$. After using
$\|(H_0^M)^{-1}\|\le1/(\Delta|M|)$, the contribution of
\eqref{sc:type-zero} is bounded by
\[
 \frac{DJ}{\Delta}\,2^k\gamma_0(s)s^k
                     \prod_i\|C_i\|_1.
\]

Now suppose $u\notin p$ and $u\in M_j$. Then
$n=|N_*\setminus p|\ge1$ and $|M_j|\le n+s$. Hence
\begin{align}
 \sum_M\frac{\|\Pi_MX\Omega\|}{|M|}
 &\le\left(\sum_{t=0}^{s}\frac{\binom{s}{t}}{(n+t)^2}
          \right)^{1/2}\|X\Omega\|
 \nonumber\\
 &\le\frac{\gamma_1(s)}{|M_j|}\|X\Omega\|.
 \label{sc:type-j}
\end{align}
For the last inequality, each function
$(n+s)/(n+t)$ is nonincreasing for $n\ge1$, since $t\le s$.
Thus its square is bounded by $(s+1)^2/(t+1)^2$, and
\eqref{sc:gammas} applies. This proves the estimate for every $n$;
no finite enumeration of sector sizes is needed.

For fixed $M_j$, at most $D|M_j|$ members $p$ intersect $M_j$.
This cancels the denominator $|M_j|$ in \eqref{sc:type-j}.
Sum the other $k-1$ supports using \eqref{sc:incidence-sum},
and then sum $M_j\ni u$. For this $j$ the result is bounded by
\[
 \frac{DJ}{\Delta}\,2^k\gamma_1(s)s^{k-1}
                         \prod_i\|C_i\|_1.
\]
Adding the $k$ possible choices of $j$ proves
\eqref{sc:multilinear}. For $k=0$, only the first case occurs,
with $N_*=\varnothing$, and its bound is $DJ\gamma_0(s)/\Delta$.
\end{proof}

\begin{proof}[Proof of Theorem~\ref{sc:stability} in finite dimension]
For a creation operator $C$, put
\begin{equation}
 W(C)=e^CVe^{-C}
     =\sum_{k=0}^{2s}\frac{1}{k!}\operatorname{ad}_C^k(V).
 \label{sc:W}
\end{equation}
The equality is exact by Lemma~\ref{sc:support}.
By \eqref{sc:free-commutator}, $e^{-C}\Omega$ is an eigenvector
of $H(\varepsilon)$ precisely when
\begin{equation}
 c_M=\varepsilon(H_0^M)^{-1}\Pi_M W(C)\Omega
 \quad(M\ne\varnothing).
 \label{sc:KT}
\end{equation}
Its eigenvalue is $E(\varepsilon)=\varepsilon\langle\Omega,W(C)\Omega\rangle$.
Indeed, conjugating $H$ by $e^C$ gives
$H_0+[C,H_0]+\varepsilon W(C)$; its nonvacuum components on
$\Omega$ vanish exactly when \eqref{sc:KT} holds.

Let $F_\varepsilon$ denote the right side of \eqref{sc:KT},
assembled as a creation operator. Lemma~\ref{sc:sector-estimate} gives
\begin{equation}
 \|F_\varepsilon(C)\|_1\le\lambda q_s(\|C\|_1).
 \label{sc:F-bound}
\end{equation}
The same multilinear estimate, applied to the difference of two
$k$-fold commutators, gives, whenever $\|C\|_1,\|C'\|_1\le r$,
\begin{equation}
 \|F_\varepsilon(C)-F_\varepsilon(C')\|_1
 \le\lambda q_s'(r)\|C-C'\|_1.
 \label{sc:contraction}
\end{equation}
For example, the difference of the $k$-th terms is a sum of $k$
multilinear terms containing one $C-C'$ and $k-1$ factors bounded
by $r$. This explains the derivative of $q_s$ in the estimate.
By \eqref{sc:criterion}, the closed ball of radius $r$ is preserved
and the contraction factor is less than $1/2$. The Banach fixed-point
theorem supplies a solution $C$ with $\|C\|_1\le r$. In finite
volume the space of creation coefficients is complete in
\eqref{sc:creation-norm}. The vector $e^{-C}\Omega$ is nonzero,
since its vacuum coefficient is one.

We next prove isolation and simplicity of this eigenvalue. Suppose
$H\phi=(E+\delta)\phi$ with $|\delta|<\Delta/2$, where $\phi$
is linearly independent of $e^{-C}\Omega$. There is a unique
creation operator $B$ and scalar $b_0$ such that
\[
 e^C\phi=(B+b_0)\Omega.
\]
Here $B\ne0$. The nonvacuum components of the eigenvalue equation
after conjugation yield
\begin{equation}
 b_M=\varepsilon(H_0^M-\delta)^{-1}
                     \Pi_M[B,W(C)]\Omega.
 \label{sc:excited-equation}
\end{equation}
To verify the sign, note that
$[H_0+[C,H_0]+\varepsilon W(C),B]
 =[H_0,B]-\varepsilon[B,W(C)]$.
The creation part $[C,H_0]$ commutes with $B$.
Since $|M|\ge1$,
\begin{equation}
 \|(H_0^M-\delta)^{-1}\|
 \le\frac{1}{\Delta|M|-|\delta|}
 \le\frac{2}{\Delta|M|}.
 \label{sc:shifted-resolvent}
\end{equation}
The proof of Lemma~\ref{sc:sector-estimate} therefore applies with
an extra factor two. In \eqref{sc:W}, commutation with $B$ increases
the order by one; orders greater than $2s$ vanish. Consequently
\begin{align}
 \|B\|_1
 &\le2\lambda\sum_{k=0}^{2s-1}
       \frac{\beta_{k+1}(s)}{k!}\|C\|_1^k\|B\|_1
 \nonumber\\
 &\le2\lambda q_s'(r)\|B\|_1<\|B\|_1,
 \label{sc:gap-contradiction}
\end{align}
a contradiction. This includes $\delta=0$, so $E$ is simple.

Finally vary the real parameter from $0$ to $\varepsilon$. The same
ball and contraction factor work throughout this segment, and
\eqref{sc:contraction} implies continuous dependence of the fixed
point on the parameter. Thus $E$ is a continuous simple eigenvalue
branch, separated by at least $\Delta/2$ from every other eigenvalue.
At zero it is the ground eigenvalue of $H_0$. Continuity of the
finite-dimensional eigenvalues prevents a crossing, proving
\eqref{sc:gap}.
\end{proof}

\begin{proof}[Extension to compact-resolvent site operators]
Let $P_{u,R}=\boldsymbol{1}_{[0,R]}(h_u)$ and
$\mathcal F_R=\bigotimes_u\Ran P_{u,R}$. These are finite-dimensional
spaces containing the product vacuum. Compress $H$ to $\mathcal F_R$.
Each onsite gap remains at least $\Delta$, and each compressed
interaction is supported on the same $p$ and has norm at most $J$.
The finite-dimensional theorem therefore applies with the same constants.

In a finite system $V$ is bounded. Hence $H$ is self-adjoint on
$\Dom H_0$, bounded below, and has compact resolvent. The nested
union of the spaces $\mathcal F_R$ contains all finite linear
combinations of product eigenvectors of the $h_u$, and is an operator
core, hence also a form core, for $H_0$ and $H$.
The min--max principle gives, for each fixed eigenvalue index $j$,
\[
 E_j(H|_{\mathcal F_R})\mathbin{\downarrow}E_j(H),
\]
where $H|_{\mathcal F_R}$ denotes compression and eigenvalues are
ordered with multiplicity. If the full space has dimension at least
two, pass to the limit in
$E_1(H|_{\mathcal F_R})-E_0(H|_{\mathcal F_R})\ge\Delta/2$.
The resulting positive difference proves both the asserted gap and
simplicity. For a one-dimensional full space, the spectral exclusion
is vacuous and the ground state is simple. No norm of an unbounded
onsite operator is used in the estimates.
\end{proof}

\subsection{Explicit rational certificates}
\label{sc:certificates}

The following choices make the criterion explicit:
\begin{equation}
 \lambda_3=0.01790402,\qquad \lambda_4=0.01190829.
 \label{sc:certified-lambdas}
\end{equation}
Every terminating decimal in this subsection denotes an exact rational
number. To check the inequalities without locating a polynomial root,
use the following rational upper bounds and witnesses:
\begin{center}
\begin{tabular}{c c c c}
$s$ & $\overline\gamma_0$ & $\overline\gamma_1$ & $r$\\
\hline
3 & 1.269295517644 & 5.859465277083 & 0.03911425\\
4 & 1.464866319271 & 8.597480251020 & 0.02996979
\end{tabular}
\end{center}
The exact squares required by \eqref{sc:gammas} are
\begin{equation}
 \begin{array}{c|cc}
 s&\gamma_0(s)^2&\gamma_1(s)^2\\ \hline
 3&29/18&103/3\\
 4&103/48&887/12
 \end{array}.
 \label{sc:exact-gamma-squares}
\end{equation}
Squaring the displayed rational upper bounds verifies
$\overline\gamma_i^2\ge\gamma_i^2$.
Let $\overline q_s$ be \eqref{sc:majorant} with these upper bounds
substituted in \eqref{sc:betas}. Direct rational substitution gives,
for each of the two rows,
\begin{equation}
 r-\lambda_s\overline q_s(r)>\frac{3}{10^8},
 \qquad
 1-2\lambda_s\overline q_s'(r)>\frac{9}{10^7}.
 \label{sc:rational-slacks}
\end{equation}
These are finite rational inequalities, verifiable by clearing
denominators. Since $q_s$ and $q_s'$ have nonnegative coefficients
bounded by those of $\overline q_s$, they imply
\eqref{sc:criterion} for all $0\le\lambda\le\lambda_s$.
The certificate checks only this arithmetic; the universal operator
conclusion comes from the preceding proof. It does not follow from
finite matrix diagonalizations or from sampling interaction operators.

\subsection{Application to the specified $\SU(N)$ lattice Hamiltonian}

Let $d\ge2$ be the spatial dimension, $N\ge2$, $x=g^2>0$, and
$0\le b_N\le2N$. On a finite collection $E$ of positively oriented
cubic-lattice links, retain any subset $\mathcal P$ of the usual
elementary plaquettes whose four links belong to $E$, with no repetitions
of an anchored plaquette. Reverse
orientations denote inverses of the positive link variables.
At lattice spacing one, use
\begin{equation}
 H_{\rm lat}(x)=\frac{x}{2}\sum_{l\in E}C_l+
       \frac{b_N}{x}\sum_{p\in\mathcal P}(1-T_p),
 \qquad T_p=\frac{\operatorname{Re}\Tr U_p}{N}.
 \label{sc:KS-H}
\end{equation}
The link space is $L^2(\SU(N),\dd U)$ for normalized Haar measure;
$C_l$ is the positive group Casimir with
$\Tr(t_at_b)=\delta_{ab}/2$. The physical space is invariant under
the vertex gauge groups, with no external charges. These conventions
specify the magnetic normalization explicitly; compare the Hamiltonian
formulation in \cite{KS}. The bounds below concern the dimensionless
energies of \eqref{sc:KS-H}.

The least nonzero Casimir is
\begin{equation}
 C_F=\frac{N^2-1}{2N}\ge\frac34.
 \label{sc:CF}
\end{equation}
For completeness, write a nonzero dominant highest weight as
$\Lambda=\omega_i+\nu$, where $\nu$ is dominant and the roots have
squared length two. The Gram entries
$(\omega_i,\omega_j)=\min(i,j)(N-\max(i,j))/N$ are positive, so
$C_2(\Lambda)\ge C_2(\omega_i)$ for
$C_2(\Lambda)=(\Lambda,\Lambda+2\rho)/2$,
where $\rho=\sum_{j=1}^{N-1}\omega_j$.
Since $C_2(\omega_i)=i(N-i)(N+1)/(2N)\ge C_F$, and the
fundamental representation attains equality, Peter--Weyl decomposition
proves \eqref{sc:CF} on the entire link space.

Divide \eqref{sc:KS-H} by $s_0=xC_F/2$ and remove its scalar
magnetic constant. Precisely,
\begin{equation}
 \frac{H_{\rm lat}}{s_0}-\varepsilon|\mathcal P|I
   =\sum_{l\in E}\frac{C_l}{C_F}
        +\varepsilon\sum_{p\in\mathcal P}(-T_p),
 \qquad
 \varepsilon=\frac{2b_N}{x^2C_F}\le\frac{32}{3x^2}.
 \label{sc:normalization}
\end{equation}
The last inequality uses
$2b_N/C_F\le8N^2/(N^2-1)\le32/3$.
The normalized onsite gap is $\Delta=1$, $J=1$ because
$\|T_p\|\le1$, and each plaquette meets $s=4$ link sites.
Each link belongs to at most $D=2(d-1)$ elementary plaquettes.
These statements hold for any finite open subgraph, and also for
the usual periodic directed cubic lattices of side lengths at least
two, with positive-direction link factors distinguished.

\begin{corollary}[Explicit $N$-uniform finite-volume gauge gap]
\label{sc:gauge-corollary}
For the model \eqref{sc:KS-H}, let
\begin{equation}
 X_d^{\rm link}
  =\left(\frac{(32/3)\,2(d-1)}{0.01190829}\right)^{1/2}.
 \label{sc:X}
\end{equation}
If $x\ge X_d^{\rm link}$, the ground state is unique and gauge
invariant, and both on the full link space and on its physical
subspace,
\begin{equation}
 H_{\rm lat}-E_0\ge\frac{xC_F}{4}\ge\frac{3x}{16}
 \quad\hbox{on the orthogonal complement of the ground state}.
 \label{sc:gauge-gap}
\end{equation}
In three spatial dimensions the sufficient condition $x\ge59.9$
implies \eqref{sc:gauge-gap}, for every $N\ge2$ and
$0\le b_N\le2N$, uniformly over the stated finite volumes.
\end{corollary}
\begin{proof}
Equations \eqref{sc:normalization}--\eqref{sc:X} give
$\lambda=\varepsilon D\le\lambda_4$.
Theorem~\ref{sc:stability} and the certificate give a normalized
gap at least $1/2$. Restoring $s_0$ gives $xC_F/4$; adding back
the scalar constant changes no gap. All Casimirs and plaquette
traces commute with the finite product of vertex gauge groups.
The unique ground line therefore carries a continuous one-dimensional
representation of a product of $\SU(N)$ groups. Such a representation
is trivial, so the ground vector is gauge invariant. The physical
space reduces $H_{\rm lat}$ and contains its ground state;
restricting the excited-space inequality cannot lower the gap.
Finally, the exact rational inequality
$(59.9)^2(0.01190829)>(32/3)\,4$ verifies the stated threshold.
\end{proof}

For reference, upward-rounded sufficient thresholds using link sites
are $42.4$, $59.9$, and $73.4$ in spatial dimensions $2$, $3$, and $4$,
respectively. The exact values from \eqref{sc:X} are approximately
$42.326$, $59.858$, and $73.310$; these approximations are not used
as downward-rounded sufficient bounds.
Alternatively, group the $d$ outgoing positive links at each vertex
into one site. Each plaquette then meets three sites, and a vertex
site meets at most $3\binom d2$ plaquettes: for each coordinate pair
its three possible roles are the lower anchor and its two forward
neighbors. With $s=3$, $D=3d(d-1)/2$, and $\lambda_3$ from
\eqref{sc:certified-lambdas}, the sufficient threshold is
$\bigl((32/3)D/\lambda_3\bigr)^{1/2}$.
In $d=2$, this permits the slightly smaller upward-rounded value
$42.3$. Link grouping gives the better displayed bound in $d=3$.

\begin{remark}[Limits of the conclusion]
\label{sc:scope}
The uniformity established here is over finite volume and $N$ at
fixed $x$. The proof does not construct an infinite-volume GNS
Hamiltonian at the explicit threshold, nor does it prove a continuum
limit or a positive physical continuum mass. An infinite-volume
strong-coupling construction with a non-explicit threshold is a
separate consequence of the stability theorem in \cite{Yarotsky},
as applied in \cite{Spectral}. A limit with fixed 't~Hooft coupling
$Nx$ eventually leaves the present window as $N$ increases.
\end{remark}

\subsection{A counting correction in the cited BDL preprint}
\label{sc:BDL-correction}

This subsection concerns the exact version
arXiv:0707.1894v1 of \cite{BDL}, Appendix~A, p.~25, in the
paragraph preceding its equation~(69). The intermediate inference
$M_j\subset M\cup\{w\}$, and hence $|M_j|\le2|M|$, does not
follow from the support conditions stated there. The following
example also verifies that this is an operator issue, rather than
only a set-theoretic possibility.

Take three distinct qubit sites $u,v,w$, their vacuum $|000\rangle$,
and
\[
 A=a_u^\dagger a_v^\dagger a_w^\dagger,
 \qquad V_{vw}=X_vX_w,
 \qquad a^\dagger=|1\rangle\langle0|,
 \quad a=|0\rangle\langle1|.
\]
Here $X$ interchanges $|0\rangle$ and $|1\rangle$.
Direct application gives
\begin{equation}
 AV_{vw}|000\rangle=0,\qquad
 V_{vw}A|000\rangle=|100\rangle,\qquad
 \langle000|a_u[A,V_{vw}]|000\rangle=-1.
 \label{sc:BDL-counterexample}
\end{equation}
Thus a nonzero term has $M_j=\{u,v,w\}$ and $M=\{u\}$.
It obeys the source's support conditions but violates
$|M_j|\le2|M|$.

The valid general estimate is
\begin{equation}
 M_j\subset M\cup p,\quad |p|\le s,\quad M\ne\varnothing
 \quad\Longrightarrow\quad
 |M_j|\le|M|+s\le(s+1)|M|.
 \label{sc:corrected-count}
\end{equation}
For two-body interactions this gives $3|M|$, which is sharp by
\eqref{sc:BDL-counterexample}. It does not invalidate the conservative
constant $2^{13}$ in the source's Lemma~3. Repeating its counting
with the factor three instead of two gives, for the possible
commutator orders $1\le k\le4$, the coefficient bound
\begin{equation}
 k\cdot4\cdot3\cdot2^k2^{k-1}
        +4\cdot2^k2^k
 \le7168<8192=2^{13}.
 \label{sc:BDL-constant-survives}
\end{equation}
Higher orders vanish, and the zeroth-order estimate is smaller.
Accordingly the correction repairs this intermediate counting step
while preserving that theorem's stated conservative constant.
No priority claim is made for identifying the correction, and no
claim about revisions other than the specified source version is
needed here.
